% ============================================================
%  PLANTILLA ESQUELETO — Informe 1
%  Formulario T-22: Informe y Presentación 1
%  Subdocumentos: 1, 2, 3, 4.1, 4.2, 5 y 13
%  Compilar con: pdflatex
% ============================================================
\documentclass[11pt,a4paper,oneside]{report}

% ------------------------------------------------------------
%  Paquetes
% ------------------------------------------------------------
\usepackage[T1]{fontenc}
\usepackage[utf8]{inputenc}
\usepackage[spanish,es-tabla]{babel}
\usepackage{lmodern}
\usepackage{mathptmx}
\usepackage{graphicx}
\usepackage{xcolor}
\usepackage{geometry}
\usepackage{booktabs}
\usepackage{longtable}
\usepackage{tabularx}
\usepackage{array}
\usepackage{multirow}
\usepackage{enumitem}
\usepackage{titlesec}
\usepackage{fancyhdr}
\usepackage{hhline}
\usepackage{caption}
\usepackage[table]{xcolor}
\usepackage[hidelinks]{hyperref}

% ------------------------------------------------------------
%  Configuración del documento
% ------------------------------------------------------------
\geometry{a4paper, top=2.8cm, bottom=2.8cm, left=2.8cm, right=2.5cm}

% ------------------------------------------------------------
%  Paleta corporativa LafroX
% ------------------------------------------------------------
\definecolor{lafroxVioleta}{HTML}{4A2C7A}
\definecolor{lafroxAzul}{HTML}{1F3A5F}
\definecolor{lafroxGrato}{HTML}{5B6770}
\definecolor{lafroxClaro}{HTML}{F4F2F7}
\definecolor{lafroxAcento}{HTML}{C9A227}
\definecolor{lafroxBg}{HTML}{F7F5FA}

\hypersetup{
  colorlinks=true,
  linkcolor=lafroxAzul,
  urlcolor=lafroxAzul,
  citecolor=lafroxAzul
}

% ------------------------------------------------------------
%  Encabezados y secciones
% ------------------------------------------------------------
\titleformat{\chapter}[display]
  {\normalfont\huge\bfseries\color{lafroxVioleta}}
  {\filright\Large\color{lafroxAzul}\chaptertitlename\ \thechapter}
  {8pt}
  {\filright}
  [\vspace{2pt}{\color{lafroxAcento}\rule{\textwidth}{1pt}}]

\titleformat{\section}
  {\normalfont\Large\bfseries\color{lafroxAzul}}
  {\thesection}{0.8em}{}
\titleformat{\subsection}
  {\normalfont\large\bfseries\color{lafroxGrato}}
  {\thesubsection}{0.8em}{}
\titleformat{\subsubsection}
  {\normalfont\normalsize\bfseries\color{lafroxGrato}}
  {\thesubsubsection}{0.8em}{}

\titlespacing*{\chapter}{0pt}{-30pt}{20pt}
\titlespacing*{\section}{0pt}{16pt}{8pt}
\titlespacing*{\subsection}{0pt}{12pt}{6pt}

% ------------------------------------------------------------
%  Encabezado y pie de página
% ------------------------------------------------------------
\pagestyle{fancy}
\fancyhf{}
\fancyhead[L]{\small\color{lafroxGrato}\textbf{LafroX} -- Informe 1}
\fancyhead[R]{\small\color{lafroxGrato}\thepage}
\renewcommand{\headrulewidth}{0.6pt}
\renewcommand{\headrule}{\hbox to\headwidth{\color{lafroxVioleta}\leaders\hrule height \headrulewidth\hfill}}
\fancypagestyle{plain}{%
  \fancyhf{}
  \fancyfoot[C]{\small\color{lafroxGrato}\thepage}
  \renewcommand{\headrulewidth}{0pt}
}

% ------------------------------------------------------------
%  Opciones de tablas
% ------------------------------------------------------------
\captionsetup{labelfont={bf,color=lafroxAzul},font=small,skip=6pt}

\newcolumntype{L}[1]{>{\raggedright\arraybackslash}p{#1}}
\newcolumntype{C}[1]{>{\centering\arraybackslash}p{#1}}

\setlist[itemize]{itemsep=2pt, topsep=4pt}
\setlist[enumerate]{itemsep=2pt, topsep=4pt}

% Comando para marcador pendiente
\newcommand{\pendiente}{{\color{lafroxAcento}\textbf{[PENDIENTE]}}}

% Comando para nota de orientación (instrucciones dentro del documento)
\newcommand{\nota}[1]{%
  \par\vspace{4pt}%
  {\small\color{lafroxGrato}\textit{#1}}%
  \par\vspace{6pt}%
}

% Comando para bloque de instrucciones
\begin{document}

% ------------------------------------------------------------
%  Portada
% ------------------------------------------------------------
\begin{titlepage}
  \thispagestyle{empty}
  \noindent
  \begin{minipage}[t][\textheight][t]{\textwidth}
    \centering
    \vspace*{1.2cm}
    {\color{lafroxVioleta}\rule{\textwidth}{2pt}}\\[0.4cm]
    {\color{lafroxVioleta}\rule{\textwidth}{0.6pt}}\\[1.2cm]

    {\Huge\bfseries\color{lafroxVioleta} LAFROX}\\[0.5cm]
    {\Large\color{lafroxAzul} Propuesta de Solución --- Distribuidora Puelche S.A.}\\[0.3cm]
    {\large\color{lafroxGrato} Licitación N° TFEP-01/2026 --- Caso 02 Logística}\\[1.0cm]

    {\color{lafroxVioleta}\rule{\textwidth}{0.6pt}}\\[0.6cm]

    {\LARGE\bfseries\color{lafroxVioleta} INFORME 1}\\[0.3cm]
    {\Large\color{lafroxAzul} Presentación Preparatoria}\\[0.2cm]
    {\normalsize\color{lafroxGrato} Formulario T-22 --- Subdocumentos 1, 2, 3, 4.1, 4.2, 5 y 13}\\[1.5cm]

    {\color{lafroxVioleta}\rule{\textwidth}{0.6pt}}\\[1.2cm]

    {\normalsize\color{lafroxGrato}
    \begin{tabular}{ll}
      \textbf{Asignatura:} & Taller de Formulación de Proyectos Informáticos (ICI-5444) \\
      \textbf{Unidad:} & Escuela de Informática, PUCV \\
      \textbf{Profesor:} & Antonio Moya Villegas \\
      \textbf{Fecha de entrega:} & 07 de septiembre de 2026 \\
    \end{tabular}
    }\\[1.2cm]

    {\color{lafroxVioleta}\rule{\textwidth}{2pt}}
  \end{minipage}
\end{titlepage}

% ------------------------------------------------------------
%  Índice general
% ------------------------------------------------------------
\tableofcontents
\newpage

% ============================================================
%  CAPÍTULO 1 — PRESENTACIÓN DE LA EMPRESA (Subdoc. 1)
% ============================================================
% ============================================================
%  CAPÍTULO 1 — PRESENTACIÓN DE LA EMPRESA (Subdoc. 1)
%  Contenido: empresa, trayectoria, innovación, equipo
% ============================================================
%  FUENTES:
%  - Subdoc. 1 (Bases_Administrativas.md §1532-1537):
%      §1534 reseña/trayectoria/capacidades/lineas de negocio.
%      §1535 estructura org., dotación, certificaciones y alianzas.
%      §1536 experiencia en la industria del caso y complejidad equivalente.
%      §1537 gobierno interno de calidad, seguridad y gestión del conocimiento.
%  - T-22, Informe 1 (§1838): presentación de la empresa.
%  - Subdoc. 12 (§1627): estructura del equipo / gobernanza (ref. transversal).
%  Ver FUENTES.md para el detalle completo.
\chapter{Presentación de la Empresa (Subdoc.\ 1)}

\nota{Este capítulo corresponde al subdoc.\ 1 del Formulario T-22. Aquí se presenta la empresa proponente: quiénes son, qué hacen, qué experiencia tienen y por qué están en condiciones de responder a esta licitación. No se describe el problema del caso ni la solución; eso va en los capítulos 2 y 3. El evaluador busca ver que la empresa es real, que tiene trayectoria relevante y que cuenta con el equipo adecuado.}

% -----------------------------------------------------------
\section{Identificación y Perfil Corporativo}
% -----------------------------------------------------------

\nota{Datos básicos de la empresa: razón social, RUT, giro, dirección, representante legal, fecha de constitución, tamaño según SERVEL (micro/pequeña/mediana/grande). Si la empresa es nueva o no tiene licitaciones previas, explicitarlo sin esconderlo. Si hay wording del caso que describes como empresa existente, ajustar la narrativa a la empresa real.}

\pendiente

% -----------------------------------------------------------
\section{Filosofía Estratégica: Misión, Visión y Sostenibilidad}
% -----------------------------------------------------------

\nota{Redactar misión, visión y compromiso ESG (ambiental, social, gobernanza). Si la empresa no tiene declaraciones formales, formularlas de forma creíble alineadas al proyecto. El evaluador busca ver coherencia entre lo que la empresa dice ser y lo que propone hacer.}

\subsection{Misión}
\pendiente

\subsection{Visión}
\pendiente

\subsection{Compromiso de Sostenibilidad y Gobernanza (ESG)}
\pendiente

% -----------------------------------------------------------
\section{Trayectoria y Casos de Éxito}
% -----------------------------------------------------------

\nota{Describir proyectos anteriores relevantes a esta licitación. Idealmente 2 a 4 casos. Para cada uno: nombre del cliente, alcance del proyecto, tecnologías usadas, resultados medibles (porcentajes, tiempos, ahorros). Si no hay casos formales, usar proyectos académicos o personales, pero ser transparente. Lo que el evaluador evalúa aquí es madurez técnica demostrada, no antigüedad de la empresa.}

\subsection{Caso de éxito 1 — Trazabilidad sanitaria}
\pendiente
\nota{Describir un caso donde la empresa haya implementado o participado en un sistema de trazabilidad de productos, idealmente con componentes de cadena de frío o cumplimiento sanitario. Si no existe, plantear un proyecto académico o de investigación relevante.}

\subsection{Caso de éxito 2 — Logística / última milla}
\pendiente
\nota{Describir un caso donde la empresa haya trabajado en optimización de rutas, gestión de flota, o distribución de consumo masivo. Puede ser un proyecto académico, una consultoría, o una contribución a software open-source relevante.}

% -----------------------------------------------------------
\section{Credenciales Organizacionales y Madurez TI}
% -----------------------------------------------------------

\nota{Certificaciones (ISO 9001, ISO 27001, SOC 2, si existen), stack tecnologico dominado por la empresa, frameworks de gestion de proyectos (Scrum, PMI), y cualquier otra credencial que respalde la capacidad tecnica. Si no hay certificaciones formales, describir la madurez tecnica del equipo: lenguajes, frameworks, cloud, bases de datos, herramientas DevOps.}

\pendiente

% -----------------------------------------------------------
\section{Capacidad Instalada para Proyectos Complejos}
% -----------------------------------------------------------

\nota{Describir la capacidad operativa: cuántos desarrolladores, infraestructura disponible, experiencia con proyectos desimilar envergadura al de esta licitación. Si la empresa es unipersonal o pequeño equipo, ser honesto y enfatizar la calidad del trabajo sobre el tamaño.}

\pendiente

% -----------------------------------------------------------
\section{Estructura del Equipo de Trabajo y Gobernanza}
% -----------------------------------------------------------

\nota{Presentar el equipo que trabajará en la propuesta y su ejecución. Nombrar roles (líder de proyecto, analista, desarrollador, tester, etc.), experiencia de cada uno y cómo se coordinarán. El evaluador busca ver que hay personas suficientes y capacitadas para el alcance propuesto. Incluir una tabla con roles, nombre y responsabilidad principal.}

\begin{tabularx}{\textwidth}{|L{3.8cm}|L{3.2cm}|X|}
\hline
\rowcolor{lafroxClaro}
\textbf{Rol} & \textbf{Nombre} & \textbf{Responsabilidad principal} \\
\hline
Jefe de Proyecto & Alex Aravena & Dirección técnica, coordinación, interfaz con el cliente \\
\hline
Arquitecto de Solución & Bastián Trejo & Diseño de arquitectura, selección tecnológica \\
\hline
Encargado de Seguridad de la Información & Álvaro Catalán & Ciberseguridad, trazabilidad de datos, cumplimiento \\
\hline
Líder de Datos & Leandro Chamorro & Modelo de datos, migración, gobernanza y analítica \\
\hline
Líder de Desarrollo & \pendiente & Implementación de componentes \\
\hline
Líder de Calidad & \pendiente & Planificación y ejecución de pruebas \\
\hline
Líder de Operación / SRE & \pendiente & Disponibilidad, observabilidad, operación \\
\hline
Líder de Implantación y Gestión del Cambio & \pendiente & Despliegue, capacitación, adopción en terreno \\
\hline
\end{tabularx}


% ============================================================
%  CAPÍTULO 2 — PROBLEMA Y NECESIDAD (Subdoc. 2)
% ============================================================
% ============================================================
%  CAPÍTULO 2 — PROBLEMA Y NECESIDAD (Subdoc. 2)
%  Contenido: contexto, necesidad, procesos AS-IS, análisis
% ============================================================
%  FUENTES:
%  - Subdoc. 2 (Bases_Administrativas.md §1539-1545) + T-22 (§1839).
%  - Narrativa del problema: íntegramente del Caso_02_Logistica.md (C),
%    Título II (proceso actual) y Título III (voces de actores).
%      Retiro sanitario: C §66, §182, §362, §428.
%      Cadena de frío/rechazo: C §250, §335, §338.
%      Conteo cíclico 2,3% / merma 1,7%: C §190, §242.
%      Envases 68.000/9.400 pallets/14%: C §244, §323.
%      OTIF 82,4% -> sobre 95%: C §301, §906.
%      Costos: C §333 ($31M), §250 ($21M), §234 ($4,2M/mes).
%  - Criterios de aceptación (retiro <2h, 100% lote, etc.): C Cap. 18 §897-918.
%  - Supuestos y 16 decisiones del 16.1: C §861+.
%  - APA 7.ª ed. exigida en AA §1545 y §1842.
%  Ver FUENTES.md para el detalle completo.
\chapter{Problema y Necesidad (Subdoc.\ 2)}

\nota{Este capítulo es el corazón de la propuesta. Aquí se demuestra que entendiste el problema de Puelche mejor que nadie. No basta resumir el caso; hay que interpretarlo, estructurarlo y revelar las consecuencias que el caso solo insinúa. El evaluador busca ver que puedes traducir narrativa en problemas concretos, medibles y priorizados. Las reglas de precedencia (Art.\ 5°) y las inconsistencias detectadas están en \texttt{AGENTS.md}.}

% -----------------------------------------------------------
\section{Contextualización del Entorno}
% -----------------------------------------------------------

\nota{Presentar el contexto macro del proyecto: la industria de distribución de consumo masivo en Chile, las tendencias de digitalización en logística, la presión regulatoria sanitaria (ISP, Ley 20.644), y la competencia. No es un resumen del caso; es la \"foto amplia\" que justifica por qué esta licitación existe.}

\subsection{El detonante: el retiro sanitario de marzo de 2026}
\nota{Narrar el evento que inicia la crisis: 14.200 llamadas, 4.800 kg recuperados, \$31M de pérdida directa, rechazo de food service por sospecha de ruptura de cadena de frío (\$21M). El caso dice que \"la estimación no era suficiente\" para la autoridad.}

\subsection{Las presiones externas que transforman la crisis en amenaza estrategica}
\nota{Identificar 3-5 presiones: competidor nacional con mejor informacion, ISP con fiscalizacion activa, food service exigiendo trazabilidad, aliados exigiendo PEDIDOS 100\% completos, rotacion de conductores (35\% en 18 meses), COSTO de no hacer nada (merma 1,7\%, OTIF 82,4\%, 4,2M/mes en diferencias de rendicion).}

\subsection{Metas convocantes de la licitación}
\nota{Citar las 4 metas del Caso Cap. 18 como objetivos que la propuesta debe cumplir: (1) retiro sanitario < 2h, 100\% del lote; (2) registro continuo de temperatura; (3) OTIF con meta; (4) preventa con stock y crédito; (5) cero pedidos perdidos/duplicados; (6) ruta automática < 20 min.}

% -----------------------------------------------------------
\section{Marco Legal y Regulatorio}
% -----------------------------------------------------------

\nota{Enmarcar el problema en la legislación chilena e internacional. No es un capítulo de derecho; es un capítulo que demuestra que la propuesta no dejará al cliente expuesto a multas o cierres. Citar leyes concretas y su impacto en el proyecto.}

\subsection{Marco legal nacional}
\nota{Ley 20.644 (trazabilidad), Ley 18.256 y su reglamento (SAG, cadena de frio), Ley 20.388 (reclamos del consumidor), DL 3.500 (FONASA/ISAPRE si aplica a conductores), Ley 20.900 (proteccion de datos personales y ciberseguridad), Ley 21.526 (fintech si hay pagos), Ley 20.644 (Ley de Transito, normativa de transporte de carga).}

\subsection{Estándares internacionales y marcos de referencia}
\nota{GS1 (estándares de identificación y trazabilidad), ISO 22000/ISO 28000 (seguridad en la cadena de suministro), GAMP 5 si aplica, frameworks de referencia como ITIL o COBIT para gestión de TI. Mencionar solo los que sean relevantes y que la propuesta efectivamente use.}

% -----------------------------------------------------------
\section{Modelo Operativo Actual (Procesos AS-IS)}
% -----------------------------------------------------------

\nota{Este es el subtítulo más extenso del capítulo. Describir cada proceso actual de Puelche tal como opera HOY, tal como lo describe el caso. No es un resumen; es una descripción técnica del flujo de información y materia. Usar el vocabulario del caso. Incluir diagramas de flujo AS-IS (formato Mermaid o PlantUML) para los procesos críticos.}

\subsection{Abastecimiento y recepción de mercadería}
\nota{Flujo desde que el bodeguero recibe mercadería hasta que queda ingresada en el ERP. El caso describe: 350SKUs, planillas impresas, contingencia manual, lote sin rastrear, vencimiento sin alerta. Incluir un diagrama de flujo AS-IS de la recepción.}

\subsection{Almacenamiento e inventario}
\nota{Distribución por temperatura (ambiente 65\%, refrigerados 20\%, congelados 15\%), conteo cíclico nocturno sobre 3.400 posiciones/mes con planillas, ajuste al cierre del mes sin investigación de causa, diferencia promedio 2,3\% del valor contado. Merma por vencimiento 1,7\%.}

\subsection{Preventa y toma de pedidos}
\nota{62 preventistas, 80\% con Windows Mobile y Pocket PC 2003, app de 2009 de proveedor desaparecido, sin conexión 6-8 horas/día, respuesta en 24-72h. 25\% con cuadernos. Stock=20\% del total, crédito manual para el 75\%.}

\subsection{Planificación de rutas}
\nota{Basada en memoria del jefe de despacho, Google Maps manual, 1.200 despachos/día, stockout 2,3\%. Optimización por conocimiento tácito, no por herramienta.}

\subsection{Preparación de pedidos y carga}
\nota{Ciclo nocturno 22:00-06:00, recepción 06:00-11:00, preparación 15:00-23:00. 350 SKUs, 120 pedidos por ciclo nocturno. Salida 05:00. Incluir diagrama de flujo del ciclo operativo diario.}

\subsection{Transporte, entrega y cobro}
\nota{200 camiones, 4.200 despachos/día, rutas de 12-14 horas, chofer-cobrador con \$180M en efectivo, afirmación de entrega por papel, rendición al día siguiente, diferencias promedio \$4,2M/mes sin investigación.}

\subsection{Devoluciones, mermas y activos retornables}
\nota{14\% de pérdida anual de envases retornables (68.000 canastillos + 9.400 pallets), devolución sin registro estructurado, conductores que autorizan descuentos por su cuenta.}

\subsection{Síntesis de ineficiencias estructurales del modelo AS-IS}
\nota{Resumir en una tabla o lista las ineficiencias críticas: (1) planillas como sistema de registro, (2) información que viaja en papel, (3) costos de servir no medidos, (4) retiro sanitario que toma días, (5) control de inventario por diferencias al cierre, (6) prevención sin stock/credito en tiempo real, (7) efectivo sin control, (8) envases sin rastrear.}

% -----------------------------------------------------------
\section{Diagnóstico de Ineficiencias}
% -----------------------------------------------------------

\nota{Estructurar el diagnóstico como una tabla que cruce: proceso afectado, tipo de ineficiencia (tecnológica, organizacional, de datos), impacto cuantificado cuando sea posible, y urgencia (crítica/alta/media/baja). El evaluador busca ver que puedes priorizar problemas, no solo listarlos.}

\begin{tabularx}{\textwidth}{|L{2.8cm}|L{3.5cm}|L{2.5cm}|X|}
\hline
\rowcolor{lafroxClaro}
\textbf{Proceso} & \textbf{Ineficiencia} & \textbf{Impacto} & \textbf{Urgencia} \\
\hline
Retiro sanitario & Sin trazabilidad lote-a-lote & 21M+ pérdida, cierre potencial & Crítica \\
\hline
Preventa & Sin stock/credito en tiempo real & Pedidos incompletos, 82,4\% OTIF & Crítica \\
\hline
Cobro en ruta & Efectivo sin control & 4,2M/mes sin investigación & Alta \\
\hline
Inventario & Planillas, ajuste al cierre & 2,3\% diferencia, 1,7\% merma & Alta \\
\hline
Envases & Sin rastreo individual & 14\% pérdida anual & Alta \\
\hline
\end{tabularx}

% -----------------------------------------------------------
\section{Análisis de Impacto}
% -----------------------------------------------------------

\subsection{Impacto financiero directo}
\nota{Sumar las pérdidas cuantificables del caso: \$31M (retiro sanitario) + \$21M (rechazo cadena de frío) + \$4,2M/mes (diferencias de rendición) + 14\% de envases + 1,7\% de mermas. Estimar el costo de no hacer nada a 12 y 24 meses.}

\subsection{Impacto operacional}
\nota{Describir las consecuencias no monetarias: demora de 3 semanas para encontrar una guía, 6-8 horas sin información del preventista, 25\% de preventistas sin registro, rotación de conductores 35\% en 18 meses, captura de conocimiento tácito del jefe de despacho.}

\subsection{Impacto regulatorio y comercial}
\nota{Consecuencias de no resolver: fiscalización del ISP, pérdida de clientes food service, impossibilidad de acceder a licitaciones publicas con trazabilidad insuficiente, perdida competitiva frente al competidor que sí tiene mejor información.}

% -----------------------------------------------------------
\section{Actores Afectados}
% -----------------------------------------------------------

\nota{Identificar y describir cada rol afectado por el problema, tal como aparecen en el Caso Título III: gerente de operaciones, jefe de bodega, jefe de despacho, gerente comercial, jefe de tesorería, conductores/preventistas, autoridad sanitaria, clientes food service. Describir qué le duele a cada uno.}

\begin{tabularx}{\textwidth}{|L{2.5cm}|L{4cm}|X|}
\hline
\rowcolor{lafroxClaro}
\textbf{Rol} & \textbf{Dolor principal} & \textbf{Necesidad insatisfecha} \\
\hline
Gerente de operaciones & No sabe cuántas entregas llegan completas & Visibilidad en tiempo real \\
\hline
Jefe de bodega & Conteo manual, merma sin investigación & Control automático de inventario \\
\hline
Jefe de despacho & Rutas por memoria, stockout 2,3\% & Planificación optimizada \\
\hline
Gerente comercial & Competidor tiene mejor info & Costo de servir medido \\
\hline
Jefe de tesorería & \$4,2M/mes sin explicación & Rendición automática \\
\hline
Conductores & Cuadernos, cobro en efectivo & Herramienta simple de registro \\
\hline
Clientes food service & Sin evidencia de cadena de frío & Trazabilidad certificada \\
\hline
\end{tabularx}

% -----------------------------------------------------------
\section{Registro de Supuestos del Caso}
% -----------------------------------------------------------

\nota{Lista de supuestos implícitos del caso que la propuesta asume como ciertos. Por ejemplo: (1) la empresa tiene conectividad 4G en la mayoría de las rutas; (2) los conductores y preventistas aceptarán capacitación; (3) el ERP existente tiene API o base de datos accesible; (4) el wMS de 2013 almacena datos rastreables; (5) el 15\% de las congeladas puede tener registro térmico; (6) los 200 camiones tienen espacio para dispositivos; (7) la autoridad sanitaria no exigirá fechas específicas de cumplimiento más allá de lo licitado.}
\pendiente


% ============================================================
%  CAPÍTULO 3 — ESQUEMA DE SOLUCIÓN Y ALCANCE (Subdoc. 3)
% ============================================================
% ============================================================
%  CAPÍTULO 3 — ESQUEMA DE SOLUCIÓN Y ALCANCE (Subdoc. 3)
%  Contenido: propuesta, componentes, alcance, exclusiones
% ============================================================
%  FUENTES:
%  - Subdoc. 3 (Bases_Administrativas.md §1547-1554) + T-22 (§1840).
%  - Cap. 17.1 del Caso (C §811-835): de la necesidad al requerimiento.
%      Catálogo RF/RNF, registro de supuestos, reglas de negocio, matriz: C §815-823.
%      Reglas de negocio viven DENTRO del registro 17.1 (no capítulo aparte).
%  - Criterios de aceptación del alcance (C Cap. 18 §897-918 y §931).
%  - Componentes <-> RT: C Cap. 15 (RT-03.10, RT-03.13, RT-02.01, RT-05.10, RT-06.01).
%    NOTA: validar códigos contra transversales (ver GG "Inconsistencias");
%    no inventar RT que la solución no cubra.
%  Ver FUENTES.md para el detalle completo.
\chapter{Esquema de Solución y Alcance (Subdoc.\ 3)}

\nota{Aquí se responde la pregunta central de la licitación: ¿qué vas a entregar y por qué esta combinación de componentes es la mejor para resolver los problemas del Cap. 2? El evaluador busca ver coherencia entre lo que identificaste como problema y lo que propones como solución. Cada componente del esquema debe responder a al menos un problema del Cap. 2. Este capítulo NO es una especificación técnica detallada; eso va en los Cap. 4 y 5. Aquí es la visión de negocio y alcance.}

% -----------------------------------------------------------
\section{Visión General de la Solución}
% -----------------------------------------------------------

\nota{Describir en 1 página (máximo 2) la propuesta de alto nivel: qué sistema se construye, cómo se integra con lo existente, y por qué esta arquitectura es la correcta para el contexto de Puelche. No entrar en detalles técnicos; esto es la \"elevator pitch\" del proyecto.}

\pendiente

% -----------------------------------------------------------
\section{Componentes del Esquema de Solución}
% -----------------------------------------------------------

\nota{Listar y describir cada componente/módulo del sistema propuesto. Para cada uno: qué hace, a qué problema del Cap. 2 responde, y qué requisitos RT (Cap. 15 del caso) habilita. Idealmente 6-10 componentes, cada uno en su propia subsección.}

\subsection{Componente 1: Motor de Trazabilidad Lote-a-Lote}
\nota{Sistema que permite rastrear cada lote desde recepción hasta entrega al cliente, incluyendo cadena de frío. Responde al problema del retiro sanitario de marzo 2026 y al Art. 7.1 de las Bases Técnicas Transversales. Debe habilitar: retiro sanitario < 2h del 100\% del lote, registro continuo de temperatura, y trazabilidad completa.}

\pendiente

\subsection{Componente 2: Gestión Inteligente de Inventario}
\nota{Sistema que reemplaza el conteo cíclico manual por control automático. Responde al problema de las 3.400 posiciones/mes en planillas, la diferencia promedio del 2,3\% y la merma del 1,7\%. Debe integrarse con el ERP existente y el WMS de 2013.}

\pendiente

\subsection{Componente 3: Preventa Móvil con Stock y Crédito en Tiempo Real}
\nota{Aplicación móvil para los 62 preventistas que reemplaza el Windows Mobile/Pocket PC 2003 y los cuadernos. Debe funcionar sin conexión (mínimo un turno completo de 14 horas), sincronizar al reconectar, y mostrar stock y crédito en tiempo real.}

\pendiente

\subsection{Componente 4: Planificación Automática de Rutas}
\nota{Sistema que reemplaza la planificación manual del jefe de despacho. Debe generar rutas optimizadas en < 20 minutos (RT-02.01, RT-02.04), considerando capacidad del vehículo, ventana horaria del cliente, y condiciones de tránsito.}

\pendiente

\subsection{Componente 5: Control de Rendición y Cobro en Ruta}
\nota{Sistema que elimina las diferencias de \$4,2M/mes en la rendición de efectivo. Debe registrar automáticamente cobros, descuentos y devoluciones, y cerrar la rendición al final del turno.}

\pendiente

\subsection{Componente 6: Panel de Costo de Servir}
\nota{Dashboard que calcula y muestra el costo real de servir a cada cliente, por zona, por ruta. Responde a la necesidad del gerente comercial de \"saber cuánto cuesta cada entrega\".}

\pendiente

\subsection{Componente 7: Gestión de Envases Retornables}
\nota{Sistema que controla los 68.000 canastillos y 9.400 pallets, reduciendo la pérdida del 14\% anual. Debe permitir rastreo por unidad o por saldo por cliente.}

\pendiente

% -----------------------------------------------------------
\section{Esquema de Integración con Sistemas Existentes}
% -----------------------------------------------------------

\nota{Diagrama de alto nivel que muestra cómo el sistema propuesto se conecta con el ERP, el WMS de 2013, y la app de preventa (migrada o reemplazada). Indicar qué datos fluyen, en qué dirección, y con qué frecuencia. La capacidad de \"no romper lo que funciona\" es clave para el evaluador.}

\pendiente

% -----------------------------------------------------------
\section{Alcance del Proyecto}
% -----------------------------------------------------------

\nota{Definir con precisión qué está incluido y qué no. El evaluador busca ver que entiendes los límites. Ser explícito sobre: (1) número de usuarios; (2) número de dispositivos; (3) número de bodegas; (4) cobertura geográfica; (5) período de operación cubierto.}

\subsection{Inclusiones}
\nota{Lista cerrada de lo que SÍ se entregará: componentes, formación, documentación, soporte post-implementación, etc.}

\pendiente

\subsection{Exclusiones}
\nota{Lista cerrada de lo que NO está incluido: reemplazo del ERP, desarrollo de hardware propio, conectividad de los clientes, etc.}

\pendiente

% -----------------------------------------------------------
\section{Requisitos Regulatorios Aplicables (RT) que Habilita}
% -----------------------------------------------------------

\nota{Tabla que cruza cada componente propuesto con los requisitos transversales (RT) del Cap. 15 del caso que habilita. Solo incluir los RT que son directamente respondidos por la solución. No inventar RT que la solución no cubre.}

\begin{tabularx}{\textwidth}{|L{3.5cm}|L{2cm}|X|}
\hline
\rowcolor{lafroxClaro}
\textbf{Componente} & \textbf{RT} & \textbf{Relación} \\
\hline
Trazabilidad lote-a-lote & RT-03.13 & Retiro sanitario < 2h \\
\hline
Gestión de inventario & RT-03.10 & Operación desconectada \\
\hline
Preventa móvil & RT-02.01 & Ruta automática < 20 min \\
\hline
Planificación de rutas & RT-02.01 & Asignación eficiente \\
\hline
Control de rendición & RT-03.12 & Replicación a nuevas unidades \\
\hline
Costo de servir & RT-05.10 & Analítica de gestión \\
\hline
Envases retornables & RT-06.01 & Control de activos \\
\hline
\end{tabularx}

% -----------------------------------------------------------
\section{Supuestos y Dependencias de la Solución}
% -----------------------------------------------------------

\nota{Listar los supuestos técnicos que deben cumplirse para que la solución funcione: conectividad 4G, API del ERP, compatibilidad del WMS, cooperación de los usuarios en capacitación, plazos de entrega de hardware, etc. Si alguno no se cumple, la propuesta se modifica.}

\pendiente


% ============================================================
%  CAPÍTULO 4 — ARQUITECTURA LÓGICA (Subdoc. 4.1)
% ============================================================
% ============================================================
%  CAPÍTULO 4 — ARQUITECTURA LÓGICA (Subdoc. 4.1)
%  Contenido: capas, módulos, tecnologías, patrones, BD
% ============================================================
%  FUENTES:
%  - Subdoc. 4 (Bases_Administrativas.md §1556-1565) + T-22 (§1841).
%      Lógica: capas, módulos, límites, interfaces (§1558).
%      Integración: servicios, contratos, mensajería, versionado (§1560).
%      Seguridad Zero Trust (§1561); decisiones registradas (§1564).
%      Prohibido diagramas genéricos (§1565).
%  - Híbrido obligatorio (AA Art. 16° y §1841).
%  - Operación desconectada 14h/24h (C §731, RT-03.10).
%  - Volumetría del caso (C Tabla 14.1: 350 SKUs, 14.200 clientes, 31.000 pedidos-mes).
%  Ver FUENTES.md para el detalle completo.
\chapter{Arquitectura Lógica (Subdoc.\ 4.1)}

\nota{Este capítulo describe CÓMO se construye la solución a nivel lógico: capas, módulos, tecnologías, patrones de diseño, bases de datos y lenguajes de programación. NO describe hardware ni infraestructura física (eso va en Cap. 5). El evaluador busca ver dominio técnico, coherencia con los componentes del Cap. 3, y que las tecnologías elegidas son apropiadas para el contexto (costo, madurez, ecosistema de talento en Chile).}

% -----------------------------------------------------------
\section{Visión General de la Arquitectura}
% -----------------------------------------------------------

\nota{Describir el estilo arquitectónico elegido (por ejemplo: arquitectura por capas, microservicios, o híbrida) y por qué es el correcto para este proyecto. Justificar la decisión: si es por capas, porque la empresa tiene equipo pequeño y necesita simplicidad; si es microservicios, porque los componentes son independientes y escalan por separado. Incluir un diagrama de alto nivel (Mermaid) que muestre las capas principales.}

\pendiente

% -----------------------------------------------------------
\section{Capas de la Arquitectura}
% -----------------------------------------------------------

\nota{Describir cada capa de la arquitectura: presentación, lógica de negocio, acceso a datos, e integración. Para cada capa: qué responsabilidad tiene, qué tecnologías la implementan, y cómo se comunica con la capa adyacente.}

\subsection{Capa de Presentación}
\nota{Interfaces de usuario: app móvil (preventistas/conductores), dashboard web (gerencia, bodega), portal de clientes (si aplica). Tecnologías: React Native o Flutter para móvil, React/Angular para web. Justificar la selección según el ecosistema de talento disponible.}

\pendiente

\subsection{Capa de Lógica de Negocio}
\nota{Motor de reglas de negocio: trazabilidad, cálculo de costo de servir, optimización de rutas, gestión de inventario, control de rendición. Tecnologías: Python (FastAPI/Django), Node.js (NestJS), o Java (Spring Boot). Justificar según el dominio y el equipo.}

\pendiente

\subsection{Capa de Acceso a Datos}
\nota{ORMs, patrones de acceso, caching. Tecnologías: SQLAlchemy (Python), TypeORM (Node), o JPA (Java). Incluir justificación de por qué se usa el patrón elegido (Active Record vs Data Mapper, etc.).}

\pendiente

\subsection{Capa de Integración}
\nota{Cómo se conecta con el ERP existente, el WMS de 2013, y servicios externos (APIs de mapas, SMS, correo). Patrones: ETL batch, webhooks, cola de mensajes. Tecnologías: Apache Kafka/RabbitMQ para eventos, REST/GraphQL para APIs síncronas.}

\pendiente

% -----------------------------------------------------------
\section{Módulos Funcionales}
% -----------------------------------------------------------

\nota{Detalle técnico de cada módulo funcional del Cap. 3: qué hace internamente, qué datos maneja, qué APIs expone, y cómo interactúa con otros módulos. Usar un formato consistente para cada módulo: nombre, responsabilidad, API pública (endpoints o eventos), dependencias, y notas de implementación.}

\subsection{Módulo de Trazabilidad}
\nota{Registrar cada evento del ciclo de vida de un lote: recepción, almacenamiento, preparación, carga, tránsito, entrega. Cada evento tiene timestamp, ubicación GPS, responsable, y condición (temperatura si aplica). API: consulta de lote por código, historial completo, alertas de excursión térmica.}

\pendiente

\subsection{Módulo de Gestión de Inventario}
\nota{Reemplazo del conteo cíclico manual: conteo por dispositivos móviles con código de barras/QR, diferencias automáticas, alertas de vencimiento, rotación ABC. Integración con el WMS existente vía API o ETL.}

\pendiente

\subsection{Módulo de Preventa Móvil}
\nota{Aplicación offline-first: catálogo de productos, stock en tiempo real, crédito del cliente, toma de pedidos, firma digital, geolocalización. Sincronización diferida cuando hay conectividad. Almacenamiento local: SQLite o Realm.}

\pendiente

\subsection{Módulo de Planificación de Rutas}
\nota{Algoritmo de optimización de rutas: vehicle routing problem (VRP) con ventanas de tiempo, capacidad del vehículo, y restricciones de zona. API de mapas: Google Directions, OSRM, o Valhalla (open source). Tiempo objetivo: < 20 minutos para 4.200 despachos.}

\pendiente

\subsection{Módulo de Rendición y Cobro}
\nota{Registro de cobros en efectivo, transferencia, cheque, descuento autorizado, devolución. Cierre automático al final del turno. Conciliación con el ERP. Alertas por diferencias > umbral.}

\pendiente

\subsection{Módulo de Dashboard de Costo de Servir}
\nota{Cálculo del costo real por entrega: combustible, tiempo de conductor, distancia, peso/volumen, zona. Dashboard con filtros por cliente, zona, período, ruta. Exportación a Excel/PDF. Integración con datos del ERP.}

\pendiente

% -----------------------------------------------------------
\section{Modelo de Datos Conceptual}
% -----------------------------------------------------------

\nota{Diagrama entitéd-relación (ER) de alto nivel que muestre las entidades principales: Lote, Producto, Cliente, Pedido, Ruta, Entrega, Temperatura, Envase, Conductor, Vehículo, Preventista, Rendición. No incluir todos los atributos; solo las relaciones clave. Usar Mermaid ER diagram. Incluir una tabla resumen de las entidades más importantes con su propósito.}

\pendiente

% -----------------------------------------------------------
\section{Tecnologías Seleccionadas}
% -----------------------------------------------------------

\nota{Tabla consolidada de todas las tecnologías usadas en la arquitectura, con justificación de cada una. Incluir: lenguaje, framework, base de datos, cola de mensajes, API de mapas, herramienta de build, plataforma de despliegue. Para cada una, justificar con al menos un criterio: madurez, costo, talento disponible en Chile, licencia, rendimiento.}

\begin{tabularx}{\textwidth}{|L{2.5cm}|L{3cm}|X|}
\hline
\rowcolor{lafroxClaro}
\textbf{Capa} & \textbf{Tecnología} & \textbf{Justificación} \\
\hline
Backend & \pendiente & \\
\hline
Frontend web & \pendiente & \\
\hline
Móvil & \pendiente & \\
\hline
Base de datos & \pendiente & \\
\hline
Cola de mensajes & \pendiente & \\
\hline
API de mapas & \pendiente & \\
\hline
Integración ERP & \pendiente & \\
\hline
\end{tabularx}

% -----------------------------------------------------------
\section{Patrones de Diseño y Buenas Prácticas}
% -----------------------------------------------------------

\nota{Describir los patrones de diseño que se usarán: Repository, Unit of Work, CQRS si aplica, Event Sourcing para trazabilidad, Clean Architecture o Hexagonal si se usa. También: versionado de APIs, manejo de errores, logging, seguridad por capas. No es un catálogo académico; es un inventario de lo que efectivamente se aplicará.}

\pendiente


% ============================================================
%  CAPÍTULO 5 — ARQUITECTURA FÍSICA (Subdoc. 4.2)
% ============================================================
% ============================================================
%  CAPÍTULO 5 — ARQUITECTURA FÍSICA (Subdoc. 4.2)
%  Contenido: infraestructura, hardware, red, nube, on-premise
% ============================================================
%  FUENTES:
%  - Subdoc. 4 (Bases_Administrativas.md §1556-1565), parte física:
%      emplazamiento nube/on-premise justificado por componente, Art. 16° (§1559).
%      despliegue: ambientes, HA, DR, respaldos (§1562).
%      dimensionamiento y plan de capacidad (§1563).
%  - Híbrido obligatorio: AA Art. 16° y TT §4.1.
%  - Emplazamiento on-premise del caso: C §738 (RT-06.01: sala ténica CD Talca
%    25 m2; gabinetes de borde CD Concepción + 3 plataformas cross-docking).
%  - Operación desconectada: C §731 (RT-03.10: CD 24h sin enlace, dispositivos 14h).
%  - Ambientes / DR: AA Art. 24°, AA Art. 20°, TT RT-04.01, RT-07.07 (ver GG #3).
%  Ver FUENTES.md para el detalle completo.
\chapter{Arquitectura Física (Subdoc.\ 4.2)}

\nota{Este capítulo describe dónde y sobre qué infraestructura vive la solución: servidores, almacenamiento, red, dispositivos móviles, componentes IoT. Las Bases Técnicas Transversales (Art. 4.1) exigen despliegue híbrido obligatorio: componente en la nube + componente on-premise. Este capítulo debe demostrar que el despliegue híbrido es viable, seguro y está correctamente dimensionado.}

% -----------------------------------------------------------
\section{Estrategia de Despliegue Híbrido}
% -----------------------------------------------------------

\nota{Justificar por qué se usa despliegue híbrido (nube + on-premise) y no solo nube o solo on-premise. En el contexto de Puelche: (1) el centro de distribución en Concepción puede perder conectividad y debe seguir operando (on-premise); (2) la analítica y el dashboard de costo de servir escalan mejor en la nube; (3) la Ley 20.900 puede exigir que ciertos datos residan en territorio nacional. Incluir un diagrama de alto nivel que muestre qué componentes van en nube y cuáles on-premise.}

\pendiente

% -----------------------------------------------------------
\section{Componente On-Premise}
% -----------------------------------------------------------

\nota{Describir la infraestructura en el centro de distribución de Concepción. Este componente debe garantizar operación continua del centro de distribución (mínimo 24 horas continuas de recepción, preparación y despacho) sin enlace de red, según RT-03.10.}

\subsection{Servidor Local (Edge)}
\nota{Especificaciones del servidor on-premise: CPU, RAM, almacenamiento (SSD/NVMe), sistema operativo, virtualización (Proxmox, VMware, o bare metal). Justificar el dimensionamiento según la carga esperada (recepciones, preparación, despacho simultáneos).}

\pendiente

\subsection{Dispositivos de Terreno}
\nota{Dispositivos para conductores y preventistas: tablets o smartphones, escáneres de código de barras/QR, dispositivos IoT para temperatura (sensores Bluetooth o LoRa). Deben operar un turno completo de 14 horas sin cobertura (RT-03.10). Incluir especificaciones mínimas y justificación.}

\pendiente

\subsection{Red Local}
\nota{Infraestructura de red dentro del centro de distribución: switches, access points WiFi, segmentación de red (VLANs), cortafuegos locales. Considerar: cobertura WiFi en bodega, carga de dispositivos, y respaldo de conectividad (4G/5G).}

\pendiente

\subsection{Alimentación Eléctrica y Continuidad}
\nota{UPS (baterías de respaldo), generador eléctrico si aplica, plan de continuidad ante cortes de energía. El centro de distribución opera 24/7 y no puede detenerse por falta de energía.}

\pendiente

% -----------------------------------------------------------
\section{Componente en la Nube}
% -----------------------------------------------------------

\nota{Describir la infraestructura en la nube (AWS, Azure, o GCP). Justificar la selección del proveedor: costo, disponibilidad en Chile (regiones cercanas), cumplimiento normativo, ecosistema de servicios.}

\subsection{Proveedor de Nube}
\nota{Comparar al menos 2 proveedores (AWS vs Azure vs GCP) y justificar la elección. Considerar: región más cercana a Chile, costo estimado mensual, servicios relevantes (compute, storage, managed databases, IoT hub).}

\pendiente

\subsection{Servicios Utilizados}
\nota{Listar los servicios específicos de nube que se usarán: EC2/VM para compute, RDS/Cloud SQL para base de datos, S3/Blob Storage para archivos, IoT Hub para sensores, CDN para contenido estático. Para cada uno: tamaño estimado, costo mensual aproximado, y justificación.}

\pendiente

\subsection{Disaster Recovery}
\nota{Plan de recuperación ante desastres: backup automático (frecuencia, retención), réplica cross-region si aplica, RTO (Recovery Time Objective) y RPO (Recovery Point Objective) definidos. Debe ser consistente con los acuerdos de nivel de servicio del Art. 4.1.}

\pendiente

% -----------------------------------------------------------
\section{Conectividad entre Componentes}
% -----------------------------------------------------------

\nota{Diagrama de red que muestre cómo se conectan los componentes: centro de distribución $\leftrightarrow$ nube $\leftrightarrow$ dispositivos móviles. Protocolos de comunicación: MQTT para IoT, HTTPS para APIs, WebSockets para actualizaciones en tiempo real. Considerar: latencia tolerable, ancho de banda, y comportamiento ante desconexión.}

\pendiente

% -----------------------------------------------------------
\section{Dimensionamiento y Costos de Infraestructura}
% -----------------------------------------------------------

\nota{Tabla consolidada del dimensionamiento de infraestructura: componentes on-premise (qty × costo unitario), servicios de nube (qty × costo mensual), dispositivos móviles (qty × costo unitario), y costo total estimado. Incluir comparación con costo actual de la operación (si los datos del caso lo permiten).}

\begin{tabularx}{\textwidth}{|L{3cm}|C{1.5cm}|L{2.5cm}|L{2.5cm}|X|}
\hline
\rowcolor{lafroxClaro}
\textbf{Componente} & \textbf{Qty} & \textbf{Costo unitario} & \textbf{Costo total} & \textbf{Justificación} \\
\hline
Servidor on-premise & 1 & \pendiente & \pendiente & Edge computing para CD \\
\hline
Tablets/dispositivos & \pendiente & \pendiente & \pendiente & Conductores y preventistas \\
\hline
Sensores temperatura & \pendiente & \pendiente & \pendiente & Cadena de frío \\
\hline
Servidor nube (mensual) & \pendiente & \pendiente & \pendiente & Analítica, dashboard \\
\hline
Base de datos nube & \pendiente & \pendiente & \pendiente & Datos transaccionales \\
\hline
\end{tabularx}

% -----------------------------------------------------------
\section{Seguridad Física y Lógica}
% -----------------------------------------------------------

\nota{Medidas de seguridad en cada capa: acceso físico al servidor on-premise (candado, registro), acceso lógico (autenticación, autorización, cifrado en tránsito y en reposo), gestión de secretos (vault), parcheo, y auditoría. Alinear con ISO 27001 si se menciona en las credenciales del Cap. 1.}

\pendiente


% ============================================================
%  CAPÍTULO 6 — MODELO Y GESTIÓN DE DATOS (Subdoc. 5)
% ============================================================
% ============================================================
%  CAPÍTULO 6 — MODELO Y GESTIÓN DE DATOS (Subdoc. 5)
%  Contenido: datos maestros, migración, analítica, gobernanza
% ============================================================
%  FUENTES:
%  - Subdoc. 5 (Bases_Administrativas.md §1567-1574) asignado al Informe 1
%    por T-22 §1844.
%      dominio de información (§1569); motor/paradigma + teorema CAP (§1570);
%      migración/saneamiento/validación (§1571); desempeño (§1572);
%      separación transaccional/analítico (§1573); calidad/retención (§1574).
%  - Retención de datos: C §734 (RT-05.10/correcto RT-16.10).
%      doc tributarios 6 años; trazabilidad vida útil+6m (mín 5 años);
%      temperatura 5 años; evidencia entrega 6 años; geolocalización 12 meses.
%  - Volumetría: C Tabla 14.1 §688-689 (~34.000 doc/año, 3.400 posiciones/mes).
%  - Planillas/cuadernos deben desaparecer: C §269.
%  Ver FUENTES.md para el detalle completo.
\chapter{Modelo y Gestión de Datos (Subdoc.\ 5)}

\nota{Este capítulo describe QUÉ datos maneja el sistema, CÓMO se almacenan, CÓMO se migran desde los sistemas actuales, y CÓMO se usa la información para generar valor (analítica). El evaluador busca ver que entiendes que los datos son el activo más valioso del proyecto y que tienes un plan concreto para migrarlos, limpiarlos, gobernarlos y explotarlos.}

% -----------------------------------------------------------
\section{Inventario de Datos del Negocio}
% -----------------------------------------------------------

\nota{Categorizar todos los datos relevantes del caso. Para cada categoría: fuente actual (ERP, WMS, planillas, cuadernos), formato, volumen estimado, calidad actual, y criticidad para el proyecto.}

\subsection{Datos Maestros}
\nota{Productos (350 SKUs), clientes (14.200), rutas, conductores (200), preventistas (62), vehículos, ubicaciones geográficas. Fuente: ERP. Calidad: variable, algunos campos incompletos.}

\pendiente

\subsection{Datos Transaccionales}
\nota{Pedidos, entregas, cobros, devoluciones, rendiciones, trazabilidad de lotes. Fuente: ERP, WMS, planillas. Volumen alto (~34.000 documentos tributarios/mes, 1.200 despachos/día). Calidad: inconsistente entre áreas.}

\pendiente

\subsection{Datos de Sensórica}
\nota{Temperatura de cadena de frío (si se implementa el IoT), ubicación GPS de vehículos, lecturas de sensores de inventario. Fuente: dispositivos nuevos. Formato: series de tiempo. Volumen: alto (lecturas cada 5-15 minutos).}

\pendiente

\subsection{Datos de Control}
\nota{Rendiciones de efectivo, descuentos autorizados, incidentes, reclamos. Fuente: cuadernos, correos, llamadas telefónicas. Calidad: pobre, dispersa, sin registro estructurado.}

\pendiente

% -----------------------------------------------------------
\section{Modelo de Datos Físico}
% -----------------------------------------------------------

\nota{Diagrama entitéd-relación detallado (no conceptual como en Cap. 4) con tablas, columnas, tipos de datos, claves primarias y foráneas. Incluir: Lote, Producto, Cliente, Pedido, DetallePedido, Entrega, Ruta, Parada, Rendicion, Cobro, Envase, SensorTemperatura, ExcursionTermica, etc. No necesariamente todas las columnas; las más relevantes.}

\pendiente

% -----------------------------------------------------------
\section{Estrategia de Migración de Datos}
% -----------------------------------------------------------

\nota{Plan concreto para migrar los datos desde el ERP, el WMS de 2013, y los registros en papel (planillas, cuadernos) hacia el nuevo sistema. Las planillas y cuadernos deben \"desaparecer como sistema de registro\" (Caso Cap. 7). La migración es un riesgo crítico; el evaluador busca ver que lo tienes planeado.}

\subsection{Fuentes y estado de calidad}
\nota{Para cada fuente: qué datos tiene, en qué formato, qué tan completos y correctos están, y qué limpieza se requiere antes de migrar. El WMS de 2013 puede tener datos en formato propietario o anticuado.}

\pendiente

\subsection{Proceso de migración ETL}
\nota{Herramienta de migración (Apache Airflow, Talend, o script personalizado), transformaciones necesarias (mapeo de códigos, estandarización de direcciones, deduplicación de clientes), validación post-migración, y plan de rollback si algo falla.}

\pendiente

\subsection{Migración de datos en papel}
\nota{Los 14.200 puntos de entrega, las rutas del jefe de despacho, y el conocimiento tácito del equipo deben digitalizarse. Proceso: digitalización de guías, captura de rutas en dispositivos GPS durante 2-4 semanas, entrevistas con conductores para validar datos de entrega.}

\pendiente

% -----------------------------------------------------------
\section{Analítica e Inteligencia de Negocio}
% -----------------------------------------------------------

\nota{Describir cómo los datos se transforman en información accionable. Incluir: dashboards operativos (retiro sanitario, inventario, OTIF), reportes de gestión (costo de servir, mermas, rendimiento de rutas), y analítica predictiva (demanda, rotura de stock, excursiones térmicas).}

\subsection{Dashboard Operativo}
\nota{Vista en tiempo real: estado de entregas del día, alertas de cadena de frío, nivel de inventario por bodega, alertas de vencimiento. Usuario: gerente de operaciones, jefe de bodega.}

\pendiente

\subsection{Dashboard de Costo de Servir}
\nota{Cálculo del costo real por cliente, por zona, por ruta. Incluye: combustible, tiempo de conductor, distance, peso/volumen, número de paradas, y costo de incidentes (rechazos, devoluciones). Usuario: gerente comercial, jefe de tesorería.}

\pendiente

\subsection{Reportes de Cumplimiento}
\nota{Reportes para la autoridad sanitaria: trazabilidad de lotes, registros de temperatura, evidencia de retiro sanitario. Formato: PDF estructurado o portal web. Frecuencia: bajo demanda o mensual.}

\pendiente

\subsection{Analítica Predictiva (futuro)}
\nota{Modelos predictivos que se pueden construir con los datos recolectados: predicción de demanda por cliente y zona, anticipación de vencimientos, detección de patrones de merma, optimización predictiva de rutas. Esto es \"futuro\" pero debe estar planteado en la arquitectura de datos.}

\pendiente

% -----------------------------------------------------------
\section{Gobernanza y Calidad de Datos}
% -----------------------------------------------------------

\nota{Políticas de calidad: quién es dueño de cada dato, cómo se validan las actualizaciones, cómo se manejan los duplicados, cuánto tiempo se retienen los datos, y cómo se protegen (cifrado, anonimización si es necesario). Alinear con Ley 20.900 (protección de datos personales) y con ISO 27001 si aplica.}

\pendiente


% ============================================================
%  CAPÍTULO 7 — INNOVACIONES (Subdoc. 13)
% ============================================================
% ============================================================
%  CAPÍTULO 7 — INNOVACIONES (Subdoc. 13)
%  Contenido: 5 innovaciones obligatorias + traza con arquitectura
% ============================================================
%  FUENTES:
%  - Art. 28° (Bases_Administrativas.md §469-481): cartera obligatoria de
%    cinco innovaciones, una por cada tipo obligatorio, sin repetir tipo.
%    Tipos: §475-481 (producto/servicio, proceso, tecnológica/arquitectura,
%    modelo de negocio/contratación, UX/sostenibilidad/impacto social).
%  - Art. 29° (§485-497): elementos exigidos por cada innovación (Formulario T-19):
%    problema concreto del caso que resuelve, tecnología que la sustenta, etc.
%    La omisión reduce la innovación a un enunciado.
%  - Subdoc. 13 (§1636-1641) + T-22 (§1842): presentación de la cartera.
%  - Pertinencia al caso y trazabilidad con arquitectura/EDT/flujo de caja (§28.2).
%  - APA 7.ª ed. exigida en tipo 3 (§479).
%  Ver FUENTES.md para el detalle completo.
\chapter{Innovaciones (Subdoc.\ 13)}

\nota{Las Bases Técnicas Transversales (Art. 3.3) exigen presentar 5 innovaciones, al menos 1 obligatoria por tipología: (1) producto o servicio, (2) proceso, (3) tecnológica o de arquitectura, (4) modelo de negocio o de contratación, (5) experiencia de usuario, sostenibilidad o impacto social. Cada innovación debe: estar alineada a las necesidades del caso, ser descrita en 2 a 4 párrafos, incluir un esquema o diagrama, y justificar por qué es innovadora (no solo una buena práctica). El evaluador busca creatividad técnica aplicada, no adjetivos vacíos.}

% -----------------------------------------------------------
\section{Innovación 1 — Producto o Servicio}
% -----------------------------------------------------------

\nota{Describir una innovación que el proponente incorpore a su propuesta de solución, que implique un cambio en el servicio entregado al cliente. Ejemplos del caso: un servicio de trazabilidad lot-a-lote con alertas automáticas para el ISP, un servicio de\"predicción de demanda por ruta\" que el cliente no pidió pero que resuelve un problema real. Debe ser concreta, medible y justificada.}

\subsection{Descripción}
\pendiente

\subsection{Por qué es innovadora}
\nota{Justificar: ¿qué tiene esta propuesta que no tiene la solución convencional del mercado? ¿Es un servicio nuevo? ¿Resuelve un problema que otros no ven? ¿Reduce costos de forma medible?}

\pendiente

\subsection{Esquema o Diagrama}
\pendiente

% -----------------------------------------------------------
\section{Innovación 2 — Proceso}
% -----------------------------------------------------------

\nota{Describir una innovación en un proceso del negocio del cliente. Ejemplos del caso: automatizar el retiro sanitario de un proceso que hoy toma días, reemplazar el conteo cíclico manual por un conteo continuo con dispositivos, digitalizar la rendición de efectivo. Debe ser un cambio real en la forma de operar, no solo una herramienta nueva sobre un proceso viejo.}

\subsection{Descripción}
\pendiente

\subsection{Por qué es innovadora}
\pendiente

\subsection{Esquema o Diagrama}
\pendiente

% -----------------------------------------------------------
\section{Innovación 3 — Tecnológica o de Arquitectura}
% -----------------------------------------------------------

\nota{Describir una innovación en la tecnología o la arquitectura del sistema. Ejemplos: uso de edge computing para operación desconectada del centro de distribución, IoT con sensores de temperatura de bajo costo, arquitectura event-driven para trazabilidad en tiempo real, uso de inteligencia artificial para optimización de rutas con restricciones reales del caso. Debe ser técnicamente justificada y viable con las tecnologías del Cap. 4.}

\subsection{Descripción}
\pendiente

\subsection{Por qué es innovadora}
\pendiente

\subsection{Esquema o Diagrama}
\pendiente

% -----------------------------------------------------------
\section{Innovación 4 — Modelo de Negocio o de Contratación}
% -----------------------------------------------------------

\nota{Describir una innovación en cómo se presta el servicio o cómo se contrata. Ejemplos: propuesta de tarifa variable por volumen procesado en vez de fija, modelo de shared services donde otros distribuidores de la zona pueden usar la plataforma (escalabilidad del negocio del proponente), propuesta de pago por resultado (porcentaje de reducción de mermas). Debe ser realista y viable económicamente.}

\subsection{Descripción}
\pendiente

\subsection{Por qué es innovadora}
\pendiente

\subsection{Esquema o Diagrama}
\pendiente

% -----------------------------------------------------------
\section{Innovación 5 — Experiencia de Usuario, Sostenibilidad o Impacto Social}
% -----------------------------------------------------------

\nota{Describir una innovación que mejore la experiencia de los usuarios finales (conductores, preventistas, bodegueros, clientes) o que tenga un impacto ambiental/social positivo. Ejemplos: app de conductores con gamificación para reducir mermas, reducción de huella de carbono por optimización de rutas, inclusión digital de preventistas que hoy usan cuadernos. Debe ser medible (encuesta de satisfacción, reducción de quejas, etc.).}

\subsection{Descripción}
\pendiente

\subsection{Por qué es innovadora}
\pendiente

\subsection{Esquema o Diagrama}
\pendiente

% -----------------------------------------------------------
\section{Traza de las Innovaciones con la Arquitectura}
% -----------------------------------------------------------

\nota{Tabla que cruce cada innovación con los componentes de la arquitectura (Cap. 4) que la habilitan, y los módulos del Cap. 3 que impacta. Esto demuestra que las innovaciones no son ideas sueltas sino que están integradas en la solución.}

\begin{tabularx}{\textwidth}{|L{2.5cm}|L{4cm}|L{4cm}|X|}
\hline
\rowcolor{lafroxClaro}
\textbf{Innovación} & \textbf{Tipología} & \textbf{Componente(s) de Arq.} & \textbf{Módulo(s) impactado(s)} \\
\hline
1 & Producto/Servicio & \pendiente & \pendiente \\
\hline
2 & Proceso & \pendiente & \pendiente \\
\hline
3 & Tecnológica & \pendiente & \pendiente \\
\hline
4 & Modelo de negocio & \pendiente & \pendiente \\
\hline
5 & UX/Sostenibilidad & \pendiente & \pendiente \\
\hline
\end{tabularx}


% ============================================================
%  ANEXO — MATRIZ DE TRAZABILIDAD INFORME 1
% ============================================================
\chapter{Anexo: Matriz de Trazabilidad del Informe 1}
\nota{Tabla que cruza cada sección del Informe 1 con el subdoc. correspondiente, los requisitos RT que aborda, y el formulario T-22 que la cubre. Entregable de coherencia interna.}

\begin{tabularx}{\textwidth}{|L{3.5cm}|C{1.5cm}|L{5cm}|X|}
\hline
\rowcolor{lafroxClaro}
\textbf{Sección del Informe} & \textbf{Subdoc.} & \textbf{RT relevantes} & \textbf{Cubierta en T-22} \\
\hline
Presentación de la empresa & 1 & --- & Sí (Informe 1) \\
\hline
Problema y necesidad & 2 & --- & Sí (Informe 1) \\
\hline
Esquema de solución y alcance & 3 & RT-03.10, RT-03.12, RT-03.13 & Sí (Informe 1) \\
\hline
Arquitectura lógica & 4.1 & RT-02.xx, RT-03.xx & Sí (Informe 1) \\
\hline
Arquitectura física & 4.2 & RT-06.xx, RT-07.xx & Sí (Informe 1) \\
\hline
Modelo de datos & 5 & RT-05.xx & Sí (Informe 1) \\
\hline
Innovaciones & 13 & --- & Sí (Informe 1) \\
\hline
\end{tabularx}

\vfill
\begin{center}
\small\color{lafroxGrato}\textit{Fin del Informe 1 --- Plantilla Esqueleto LafroX}
\end{center}

\end{document}
